\section{Methodology}

\subsection{SPH Approximation}

For a quantity $A$ evaluated at particle $i$, the SPH approximation is written
as
\begin{equation}
    A_i \approx \sum_j m_j \frac{A_j}{\rho_j} W(\lVert \mathbf{r}_i-\mathbf{r}_j\rVert,h),
\end{equation}
where $m_j$ is the particle mass, $\rho_j$ is the density, $h$ is the smoothing
length, and $W$ is a compact-support kernel \cite{lucy1977,monaghan1992}.
We use the two-dimensional cubic spline with support radius $2h$,
\begin{equation}
    W(q,h) = \frac{5}{14\pi h^2}
    \begin{cases}
        (2-q)^3 - 4(1-q)^3, & 0 \le q < 1,\\
        (2-q)^3,            & 1 \le q < 2,\\
        0,                  & q \ge 2,
    \end{cases}
    \qquad q = \frac{\lVert \mathbf{r}_i-\mathbf{r}_j\rVert}{h}.
\end{equation}

Particles are initialised on a square lattice with spacing $\Delta x$ and
carry the mass $m = \rho_0 \Delta x^2$, so that a particle inside a regular
lattice sums to the rest density $\rho_0$. The smoothing length is set to
$h = 1.3\,\Delta x$, which gives about 20 neighbours per particle. Larger ratios
make the cubic spline prone to the pairing instability \cite{dehnen2012}, in
which particles clump into rows; this is examined in
Section~\ref{sec:pairing}.

\subsection{Density and Equation of State}

The density is obtained by direct summation over all particles within the
kernel support, including the particle itself and the wall particles,
\begin{equation}
    \rho_i = m_i W(0,h) + \sum_{j \ne i} m_j W_{ij}.
\end{equation}
The fluid is treated as weakly compressible \cite{monaghan1994,becker2007}.
Pressure follows the Tait equation
\begin{equation}
    p_i = B\left[\left(\frac{\rho_i}{\rho_0}\right)^{\gamma} - 1\right],
    \qquad B = \frac{\rho_0 c^2}{\gamma}, \qquad \gamma = 7,
    \label{eq:tait}
\end{equation}
where $c$ is the numerical speed of sound; negative pressures are set to zero
to avoid tensile clumping at the free surface. The density error scales with
the squared Mach number $\mathrm{Ma}^2 = (v/c)^2$, so $c$ is chosen as ten times
the largest expected flow speed. Estimating that speed by free fall from the top
of the fluid body of height $H$ gives
\begin{equation}
    c = 10\sqrt{2gH},
\end{equation}
which targets density errors of about one percent.

\subsection{Momentum Equation}

The acceleration of a fluid particle is the sum of a pressure term, a viscous
term and gravity,
\begin{align}
    \mathbf{a}_i^{p} &= -\sum_j m_j \left(\frac{p_i}{\rho_i^2} + \frac{p_j}{\rho_j^2}\right) \nabla_i W_{ij},\\
    \mathbf{a}_i^{\nu} &= 2\nu \sum_j \frac{m_j}{\rho_j}
        \frac{\mathbf{v}_{ij}\cdot\mathbf{r}_{ij}}{\lVert\mathbf{r}_{ij}\rVert^2 + 0.01h^2}\,\nabla_i W_{ij},\\
    \mathbf{a}_i &= \mathbf{a}_i^{p} + \mathbf{a}_i^{\nu} + \mathbf{g},
\end{align}
with $\mathbf{r}_{ij} = \mathbf{r}_i - \mathbf{r}_j$ and
$\mathbf{v}_{ij} = \mathbf{v}_i - \mathbf{v}_j$. The symmetric pressure term
conserves linear momentum; the viscous term follows
\cite{monaghan1992,morris1997} and also damps the acoustic oscillations that a
weakly compressible fluid carries.

\subsection{Boundary Conditions}

Walls are built from fixed particles placed on the same lattice as the fluid.
Enough layers are used to fill the full kernel support of a fluid particle
touching the wall ($\lceil 2h/\Delta x \rceil = 3$ layers for $h = 1.3\Delta x$),
so that wall particles enter the density sum like fluid at rest density.

Wall particles carry no state of their own, so a pressure has to be assigned to
them. If the wall term used only $p_i/\rho_i^2$ \cite{akinci2012}, the wall
would provide half the support that a fluid neighbour provides; the first fluid
layer then has to compress until its pressure is roughly doubled, which pushes
the second layer away and leaves a visible gap above the floor. Instead the wall
pressure is extrapolated from the surrounding fluid following
Adami et al.~\cite{adami2012},
\begin{equation}
    p_b = \frac{\sum_f p_f W_{bf} + \mathbf{g}\cdot\sum_f \rho_f\,(\mathbf{r}_b - \mathbf{r}_f)\,W_{bf}}{\sum_f W_{bf}},
\end{equation}
where the sums run over the fluid neighbours $f$ of the wall particle $b$ and the
second term adds the hydrostatic pressure difference across the wall. The wall
density $\rho_b$ follows from inverting Eq.~\eqref{eq:tait}, and the fluid--wall
pressure force uses the same symmetric form as the fluid--fluid force.

As a safety net against penetration during violent impacts, a fluid particle
that comes closer than $\Delta x/2$ to a wall particle is projected back to that
distance and its normal velocity is reflected with a restitution of $0.5$.
Fluid at rest sits a full spacing away from the wall, so the projection does not
act in quiescent regions.

\subsection{Time Integration}

Positions and velocities are advanced with the semi-implicit (symplectic) Euler
scheme,
$\mathbf{v}^{n+1} = \mathbf{v}^{n} + \Delta t\,\mathbf{a}^{n}$,
$\mathbf{r}^{n+1} = \mathbf{r}^{n} + \Delta t\,\mathbf{v}^{n+1}$.
The time step is chosen adaptively in every substep as the most restrictive of
the acoustic CFL, force, and viscous conditions \cite{monaghan1992},
\begin{equation}
    \Delta t = \min\left(
        0.4\,\frac{h}{c + \max_i \lVert\mathbf{v}_i\rVert},\;
        0.25\sqrt{\frac{h}{\max_i \lVert\mathbf{a}_i\rVert}},\;
        0.125\,\frac{h^2}{\nu}
    \right).
\end{equation}
Each rendered frame advances the simulation by $1/60\,\mathrm{s}$, split into as
many substeps as these limits require. For the stiff fluids used here the
acoustic condition dominates.

\subsection{Neighbour Search}

Neighbour lists are rebuilt once per substep by testing all particle pairs
against the kernel support radius $2h$. The cost of this search is
$\mathcal{O}(N^2)$; it is adequate for the scene sizes used in this work
(about $1.6\times10^3$ particles) but is the main limitation for larger scenes.

\subsection{Software Structure}

The solver is separated into the following components:
\begin{itemize}
    \item \texttt{Particle} stores particle state and physical properties.
    \item \texttt{NeighborSearch} constructs local particle neighborhoods.
    \item \texttt{SPHSolver} performs the density, wall-pressure, force, and
          integration steps and chooses the adaptive time step.
    \item \texttt{Scenes} builds the test scenes (tank walls, fluid body, and
          matching solver parameters).
    \item \texttt{Renderer} displays the scene and exposes runtime controls.
\end{itemize}

This separation allows the numerical code to be tested and run without an
interactive graphics window. A headless program (\texttt{sph\_experiments})
runs all experiments of Section~\ref{sec:results} and writes their data as CSV
files, from which the figures are generated.

\subsection{Experimental Setup}

All quantities are given in screen units (pixels and seconds). Results are
reported in dimensionless form ($p/\rho_0 g H$, $Z/a$, $T = t\sqrt{g/a}$), so the
choice of units does not affect them. Two scenes are used
(Table~\ref{tab:setup}):
\begin{itemize}
    \item \textbf{Hydrostatic column.} An open tank of width $300\,\mathrm{px}$ is
          filled to $H = 400\,\mathrm{px}$ with fluid at rest. The column has to
          settle to the hydrostatic pressure $p = \rho_0 g d$ at depth $d$.
    \item \textbf{Dam break.} A fluid column of width $a = 200\,\mathrm{px}$ and
          height $H = 400\,\mathrm{px}$ (aspect ratio $n^2 = H/a = 2$) stands in
          the lower-left corner of a closed tank of width $800\,\mathrm{px}$ and
          collapses under gravity at $t = 0$ \cite{martin1952}.
\end{itemize}

\begin{table}[ht]
    \centering
    \caption{Default parameters of both scenes. Parameter studies vary one
    quantity at a time.}
    \label{tab:setup}
    \begin{tabular}{lll}
        \toprule
        Quantity & Symbol & Value \\
        \midrule
        Particle spacing & $\Delta x$ & $10\,\mathrm{px}$ \\
        Smoothing length & $h$ & $13\,\mathrm{px}$ ($1.3\,\Delta x$) \\
        Rest density & $\rho_0$ & $300$ \\
        Gravity & $g$ & $200\,\mathrm{px/s^2}$ \\
        Speed of sound & $c$ & $10\sqrt{2gH} = 4000\,\mathrm{px/s}$ \\
        Viscosity & $\nu$ & $150\,\mathrm{px^2/s}$ \\
        Fluid / wall particles (column) & & $1200$ / $420$ \\
        Fluid / wall particles (dam break) & & $800$ / $828$ \\
        Frame time & & $1/60\,\mathrm{s}$ \\
        \bottomrule
    \end{tabular}
\end{table}
